\documentclass[12pt]{article}
\usepackage[amsmath]{e-jc}
\usepackage{booktabs}
\usepackage{longtable}

\MSC{05A18, 05C45, 05C85}

% Initial-submission presentation: retain the journal's typography without
% displaying unassigned volume, paper number, DOI, or publication dates.
\makeatletter
\renewcommand{\maketitle}{\date{}\old@maketitle}
\renewcommand{\ps@plain}{%
  \renewcommand{\@oddfoot}{\hfil\thepage\hfil}%
  \renewcommand{\@evenfoot}{\hfil\thepage\hfil}}
\pagestyle{plain}
\makeatother

\ifdefined\publicmode
  \newcommand{\blindmode}{0}
\else
  \newcommand{\blindmode}{1}
\fi

\title{Hamilton Cycles and Paths in the Noncrossing Partition Refinement Graph}

\ifnum\blindmode=1
  \author{Anonymous}
  \authortext{}{Anonymous Institution.}
\else
  \author{Alex Chengyu Li}
  \authortext{}{Independent Researcher, Tokyo, Japan (\email{cl7076@nyu.edu}).}
\fi

\newcommand{\NCR}{\mathrm{NC}_R}
\newcommand{\NC}{\mathrm{NC}}
\newcommand{\catalan}{\mathrm{Cat}}
\newcommand{\leanname}[1]{\nolinkurl{#1}}

\begin{document}
\maketitle

\begin{abstract}
Let $\NCR(n)$ be the cover graph of the refinement order on the
noncrossing partitions of an $n$-element cyclically ordered set.  Thus an
edge merges two blocks, or reverses such a merge by splitting one block.
We determine exactly when this graph has a Hamilton cycle and when it has a
Hamilton path.  If a one-vertex graph is regarded as Hamiltonian via its
constant closed spanning walk, the cycle classification for all natural $n$
is
\[
 \NCR(n)\text{ is Hamiltonian}
 \quad\Longleftrightarrow\quad
 n\in\{0,1\}\ \text{or}\ (n\text{ is even and }n\geq4).
\]
With the analogous convention for a one-vertex path, $\NCR(n)$ has a Hamilton
path exactly when $n\leq3$ or $n$ is even.  For odd $n=2m+1$, block-count
parity is a bipartition, and a signed Dyck-tree recurrence gives color-class
difference $(-1)^{m+1}\catalan_m$.  This rules out Hamilton cycles for odd
$n\geq3$ and Hamilton paths for odd $n\geq5$; the remaining small orders are
handled directly.  For every even $n\geq4$, we give a uniform Hamilton-cycle
construction.  A noncrossing partition is encoded bijectively by a
noncrossing partial matching $Q$ and a Boolean mask on the unmatched nonroot
points.  Each fixed $Q$ is an odd-dimensional hypercube with a cyclic
reflected Gray code.  Deleting the lexicographically first chord makes the
matchings into a rooted tree.  Explicit refinement diamonds and an injective
port assignment allow recursive square switching to join all Boolean cycles
into one Hamilton cycle.
\end{abstract}

\section{Introduction}

Write $X_n=\{0,1,\ldots,n-1\}$ with its cyclic order.  A partition of
$X_n$ is \emph{noncrossing} if there are no $a<b<c<d$ for which $a,c$
belong to one block and $b,d$ belong to another.  The set $\NC(n)$ of all
noncrossing partitions is the Kreweras lattice under refinement~\cite{K,S,SU}.
\paragraph{Refinement graph.}\label{def:refinement-graph}
Its cover graph $\NCR(n)$ has vertex set $\NC(n)$; two vertices are adjacent
when one is obtained from the other by merging exactly two blocks, with the
result still noncrossing.

\paragraph{Why these are precisely the cover edges.}
The Kreweras lattice is graded by
\[
 \rho(\pi)=n-b(\pi),
\]
where $b(\pi)$ is the number of blocks.  If $\sigma$ is obtained from $\pi$
by merging two blocks, then $\pi<\sigma$ and
$\rho(\sigma)=\rho(\pi)+1$, so no lattice element can lie strictly between
them.  Conversely, a cover raises the rank by one.  Since every block of a
coarsening is a union of blocks of the finer partition, lowering the number
of blocks by one can only replace two blocks by their union.  Thus the
merge/split adjacency used below is exactly the cover relation, rather than
a larger auxiliary graph.

The adjacency here is not the element-transfer adjacency in the cyclic Gray
code of Huemer, Hurtado, Noy, and Oma\~na-Pulido~\cite{HHNP}.  In their graph,
one element is moved between blocks in a prescribed local way; in our graph,
an edge changes the number of blocks by exactly one through a whole-block
merge or split.  The distinction is already visible at $n=3$: their graph is
Hamiltonian, whereas the five-vertex refinement cover graph is bipartite with
unequal color classes and hence is not Hamiltonian.  Thus the known cyclic
Gray code does not settle the problem considered here.  The general
permutation-language results of Hartung--Hoang--M\"utze--Williams and
Hoang--M\"utze~\cite{HHMW,HM} provide Hamilton paths for many quotient
graphs, but do not identify the refinement cover graph with such a quotient.
Merino, Namrata, and Williams~\cite{MNW} consider the algorithmic problem in
which the input is an arbitrary list of noncrossing partitions and one asks
whether the list admits a refinement Gray code; they prove that this decision
problem is NP-complete.  That arbitrary-sublist problem is different from the
spanning-cycle problem for the full lattice considered here.  Nevertheless,
our classification and signed enumeration completely resolve its canonical
full-family instances: all noncrossing partitions of $[n]$ admit a refinement
Gray code if and only if $n\leq3$ or $n$ is even, and for even $n\geq4$ the
Gray code may be chosen cyclically.  For broader background on combinatorial
Gray codes and their Hamiltonian formulations, see M\"utze's survey~\cite{M}.

Petersen~\cite{Pet} constructed a symmetric Boolean decomposition of the
noncrossing partition lattice using valley hopping.  Here the corresponding
coordinates are described as follows: a Boolean block is indexed by a canonical
noncrossing partial matching, and its Boolean coordinate is an explicit mask.
The decomposition alone does not provide a Hamilton cycle: one must choose a
different Gray edge for every child block, prove that the choices do not
collide as directed edge occurrences, and retain one additional edge for the
parent.  The construction below supplies this compatibility through three
types of ports and a recursive edge-piece assembly.

Our main result is the following complete classification.

\begin{theorem}[Complete classification]\label{thm:classification}
Regard a one-vertex graph as Hamiltonian via its constant closed spanning
walk.  Then, for every $n\in\mathbb N$,
\[
 \NCR(n)\text{ has a Hamilton cycle}
 \quad\Longleftrightarrow\quad
 n=0\ \text{or}\ n=1\ \text{or}\
 (n\text{ is even and }4\leq n).
\]
Consequently, in the usual nondegenerate range $n\geq3$, $\NCR(n)$ has a
Hamilton cycle if and only if $n$ is even.
\end{theorem}

\begin{corollary}[Hamilton paths]\label{cor:hamilton-paths}
Regard a one-vertex graph as having a Hamilton path of length zero.  Then, for
every $n\in\mathbb N$,
\[
 \NCR(n)\text{ has a Hamilton path}
 \quad\Longleftrightarrow\quad
 n\leq3\ \text{or}\ n\text{ is even}.
\]
\end{corollary}

\begin{proof}
For $n=0,1$ the graph has one vertex, and for $n=2$ it is a single edge.  For
$n=3$, there are three two-block partitions, each adjacent to both the
one-block and three-block partitions, so alternating through these five
vertices gives a Hamilton path.  If
$n\geq4$ is even, Theorem~\ref{thm:classification} gives a Hamilton cycle,
and deleting one edge gives a Hamilton path.  Finally, let $n=2m+1\geq5$.
By Theorem~\ref{thm:signed-count}, the two block-parity classes differ in
size by $\catalan_m\geq2$.  A path in a bipartite graph alternates between
the two classes, so no Hamilton path can exist.
\end{proof}

Theorem~\ref{thm:classification} is proved in Sections~\ref{sec:odd}
and~\ref{sec:coordinates}--\ref{sec:assembly}, which establish its two
directions.

The constructive proof has the following dependencies:
\[
\begin{array}{c}
\text{matching--mask bijection}
\longrightarrow \text{Gray cycles in the Boolean blocks}\\[-1pt]
\downarrow\\[-1pt]
\text{rooted matching tree}
\longrightarrow \text{explicit joining diamonds}\\[-1pt]
\downarrow\\[-1pt]
\text{injective directed ports}
\longrightarrow \text{recursive subtree cycle}\\[-1pt]
\downarrow\\[-1pt]
\text{Hamilton cycle}.
\end{array}
\]
Each arrow is justified below.

\section{The odd-order obstruction}\label{sec:odd}

Let $b(\pi)$ denote the number of blocks of $\pi$.  Every cover changes
$b(\pi)$ by one, hence block-count parity colors the graph properly.

\begin{lemma}\label{lem:bipartite}
The graph $\NCR(n)$ is bipartite, with color classes
\[
 E_n=\{\pi:b(\pi)\text{ is even}\},\qquad
 O_n=\{\pi:b(\pi)\text{ is odd}\}.
\]
\end{lemma}

\begin{proof}
A cover relation merges two blocks or splits one block into two.  Thus the
two endpoints have block counts differing by exactly one, and therefore have
opposite parity.
\end{proof}

Define the signed count
\[
 D_n:=\sum_{\pi\in\NC(n)}(-1)^{b(\pi)}=|E_n|-|O_n|.
\]
Use the standard Dyck-word, or equivalently plane-binary-tree,
correspondence for noncrossing partitions in which the number of blocks is
the number of peaks~\cite{S,SU}.  We count only nonempty tree nodes: the empty
tree has size zero, trees of size $n$ correspond to Dyck words of semilength
$n$, and hence to $\NC(n)$.  Let $p(T)$ be the peak count.  A tree with $n+1$
nodes has the unique form $T=\operatorname{node}(L,R)$ with
$|L|+|R|=n$.  If $L$ is empty, then $p(T)=1+p(R)$, so these trees contribute
$-D_n$.  If $L$ is nonempty, write $|L|=k+1$; then
$p(T)=p(L)+p(R)$ and $|R|=n-1-k$.  Summing over the possible subtree sizes
therefore gives
\begin{equation}\label{eq:signed-recurrence}
 D_{n+1}=-D_n+\sum_{k=0}^{n-1}D_{k+1}D_{n-1-k},
 \qquad D_0=1,\quad D_1=-1.
\end{equation}

\begin{theorem}[Signed count]\label{thm:signed-count}
For $m\geq1$,
\[
 D_{2m}=0,
 \qquad
 D_{2m+1}=(-1)^{m+1}\catalan_m,
\]
where $\catalan_m=\frac1{m+1}\binom{2m}{m}$.
\end{theorem}

\begin{proof}
We use simultaneous strong induction on $m$.  In
\eqref{eq:signed-recurrence} for $D_{2m}$, the two indices in every product
sum to $2m-1$.  Hence one index is positive and even, except for the boundary
term $D_{2m-1}D_0$.  Every nonboundary term vanishes by induction, and the
boundary term cancels the initial $-D_{2m-1}$.  Thus $D_{2m}=0$.

For $D_{2m+1}$, all terms having a positive even index vanish.  The remaining
terms have indices $2j+1$ and $2(m-1-j)+1$.  By induction their product is
\[
 (-1)^{j+1}\catalan_j
 (-1)^{m-j}\catalan_{m-1-j}
 =(-1)^{m+1}\catalan_j\catalan_{m-1-j}.
\]
The initial term is zero because $D_{2m}=0$.  Summing and using the Catalan
recurrence
$\catalan_m=\sum_{j=0}^{m-1}\catalan_j\catalan_{m-1-j}$ proves the formula.
The initial values $D_0=1$ and $D_1=-1$ start the induction.
\end{proof}

\begin{theorem}[Odd-order obstruction]\label{thm:odd-obstruction}
If $n\geq3$ is odd, then $\NCR(n)$ has no Hamilton cycle.
\end{theorem}

\begin{proof}
Write $n=2m+1$.  Since $n\geq3$, $m\geq1$, and
Theorem~\ref{thm:signed-count} gives
$|E_n|-|O_n|=(-1)^{m+1}\catalan_m\neq0$.  Every cycle in a bipartite graph
uses equally many vertices from its two color classes.  A spanning cycle
would therefore force $|E_n|=|O_n|$, a contradiction.
\end{proof}

\section{Canonical matching--mask coordinates}\label{sec:coordinates}

We now assume $n>0$.  A \emph{canonical matching} is a set $Q$ of ordered
chords $(a,b)$ with
\[
 0<a<b<n,
\]
such that distinct chords have disjoint endpoints and no two chords cross.
The chords may be nested.  Let $\operatorname{end}(Q)$ be their endpoint set
and define the free points
\[
 F(Q)=X_n\setminus(\{0\}\cup\operatorname{end}(Q)).
\]
For a chord endpoint $x$, let $e_Q(x)$ be the left endpoint of its chord.
For $x\in F(Q)$, let $a_Q(x)$ be the left endpoint of the innermost chord
whose open interval contains $x$, or $0$ if there is no such chord.

For a mask $A\subseteq F(Q)$ define
\begin{equation}\label{eq:decode-key}
 \kappa_{Q,A}(x)=
 \begin{cases}
 e_Q(x),&x\in\operatorname{end}(Q),\\
 a_Q(x),&x\in A,\\
 x,&\text{otherwise}.
 \end{cases}
\end{equation}
Let $P_Q(A)$ be the partition into the fibers of $\kappa_{Q,A}$.

\begin{lemma}[Key-fiber geometry]\label{lem:key-fibers}
If two distinct points have the same decoded key, then that key is either
$0$ or the left endpoint of a chord of $Q$.  A fiber based at a chord
$(a,b)$ is contained in $[a,b]$.  A point in that fiber cannot lie strictly
inside a properly nested chord with a different left endpoint, and a point
in the root fiber cannot lie strictly inside any chord of $Q$.
Consequently, two distinct decoded fibers cannot cross.
\end{lemma}

\begin{proof}
The first assertion follows directly from~\eqref{eq:decode-key}.  A point can
change its own key only by being a chord endpoint or by being selected in the
mask.  In the first case the common key is that chord's left endpoint.  In
the second it is the left endpoint of the innermost containing chord, or
$0$.

Fix $(a,b)\in Q$.  Its endpoints decode to $a$.  If another point $x$ also
decodes to $a$, then either $x$ is one of these endpoints or $x$ is selected
and $(a,b)$ is its innermost containing chord; in either case $a\leq x\leq b$.
If $(c,d)$ is properly nested in $(a,b)$ and $c<x<d$, then $(c,d)$ is a
strictly inner containing chord, so the anchor of a selected $x$ is at least
$c>a$; hence $x$ cannot decode to $a$.  The same argument with no outer chord
shows that a root-fiber point cannot lie inside any chord.

Suppose fibers based at $r$ and $s$ crossed, with $x_1<y_1<x_2<y_2$ in the
two fibers.  If one key is $0$, a point of the root fiber lies strictly
between two points of the chord fiber and hence strictly inside its base
chord, contradicting the root assertion.  Suppose now that the two fibers are
based at distinct chords $e=(a,b)$ and $f=(c,d)$.  Since $Q$ is noncrossing
and endpoint-disjoint, their intervals are either disjoint or one is properly
nested in the other.  Disjoint intervals cannot contain alternating points.
If $f$ is nested in $e$, then the $e$-fiber point $x_2$, which lies between
$y_1$ and $y_2$, lies strictly inside $f$, contradicting the nested-chord
assertion for the $e$-fiber.  If $e$ is nested in $f$, the same argument uses
the $f$-fiber point $y_1$ between $x_1$ and $x_2$.  These cases exhaust the
relative positions of the two base chords, so distinct fibers cannot cross.
\end{proof}

\begin{theorem}[Boolean decomposition]\label{thm:boolean-decomposition}
For every $n>0$, the map
\[
 (Q,A)\longmapsto P_Q(A),
 \qquad Q\text{ canonical},\quad A\subseteq F(Q),
\]
is a bijection onto $\NC(n)$.  Moreover, if $q\in F(Q)\setminus A$, then
$P_Q(A\cup\{q\})$ covers $P_Q(A)$ in the refinement lattice.
\end{theorem}

\begin{proof}
First note that $\kappa_{Q,A}(x)\leq x$.  By
Lemma~\ref{lem:key-fibers}, four alternating points cannot lie in two
different fibers.  Thus $P_Q(A)$ is noncrossing.

Conversely, for $\pi\in\NC(n)$, take from every nonsingleton block not based
at $0$ the chord joining its minimum and maximum.  Distinct blocks give
endpoint-disjoint noncrossing chords, so this is a canonical matching
$Q(\pi)$.  Put in $A(\pi)$ exactly those remaining nonroot, nonendpoint
points whose block is nontrivial.

We justify the anchor used by the decoder.  Let $x$ be such a remaining
point, and let $B$ be its block.  If $0\in B$, then no chord belonging to a
different block can contain $x$ in its open interval: its two endpoints,
together with $0$ and $x$, would give alternating points in two blocks.
Hence $a_{Q(\pi)}(x)=0=\min B$.  If $0\notin B$, the chord
$(\min B,\max B)$ belongs to $Q(\pi)$ and contains $x$.  Any chord properly
inside it belongs to a block nested inside $B$ and cannot contain a point of
$B$; any chord outside it is not innermost.  Thus the innermost chord
containing $x$ is the chord of $B$ itself, and
$a_{Q(\pi)}(x)=\min B$.  Chord endpoints decode to the same minimum by
definition, while singleton points retain their own key.  Therefore
\[
 \kappa_{Q(\pi),A(\pi)}(x)=\min(\text{the block of }x)
\]
for every $x$, and hence $P_{Q(\pi)}(A(\pi))=\pi$.

Conversely, start from $P_Q(A)$.  Lemma~\ref{lem:key-fibers} says that the key
of any nonroot nonsingleton fiber is the left endpoint $a$ of a chord
$(a,b)\in Q$.  Both $a$ and $b$ lie in that fiber.  Every other point with key
$a$ lies in $[a,b]$, so $a$ and $b$ are its minimum and maximum.  A second
chord cannot have the same left endpoint because chord endpoints are
disjoint.  Thus taking the extrema of all nonroot nonsingleton fibers recovers
exactly the chords of $Q$, not merely a matching with the same fibers.  Once
$Q$ is known, a free point belongs to $A$ exactly when its decoded key is not
itself, equivalently when it lies in a nonsingleton fiber.  This recovers $A$.
The two reconstruction procedures are inverse, proving both injectivity and
surjectivity.

Finally, if $q\notin A$, then $q$ is a singleton fiber because its key is
$q$.  Inserting $q$ into the mask changes only this key, from $q$ to
$a_Q(q)$.  It therefore merges the singleton $\{q\}$ with exactly the fiber
based at $a_Q(q)$ and leaves every other fiber unchanged.  The block count
drops by one, so this is one refinement cover.
\end{proof}

The dimension of the block indexed by $Q$ is
\begin{equation}\label{eq:free-ledger}
 d(Q):=|F(Q)|=n-1-2|Q|.
\end{equation}
If $n$ is even, then $d(Q)$ is positive and odd.

\section{Gray cycles, the block tree, and ports}\label{sec:ports}

Write $d=d(Q)$ and $F(Q)=\{u_0<u_1<\cdots<u_{d-1}\}$.  We use the binary-reflected Gray
list with $u_0$ as the least significant coordinate.  Thus, with
$G_0=(\varnothing)$,
\[
 G_d=G_{d-1}\mathbin{\Vert}
 \bigl(\operatorname{rev}(G_{d-1})+u_{d-1}\bigr),
\]
where $+u_{d-1}$ inserts $u_{d-1}$ into every mask.  Equivalently, the
recursive coordinate list is supplied in descending free-point order.
Adjacent masks differ by one free point, so
Theorem~\ref{thm:boolean-decomposition} sends every Gray edge to a refinement
edge.

\begin{lemma}[Block Gray cycle]\label{lem:block-gray-cycle}
For even $n$, every block $\{P_Q(A):A\subseteq F(Q)\}$ has a cyclic Gray
listing that visits each of its vertices once.  For $d(Q)\geq3$ it is a
simple cycle; for $d(Q)=1$ it is the two directed occurrences of the unique
edge, to be consumed by the tree switch below.
\end{lemma}

\begin{proof}
The displayed recursion lists every subset once, changes one coordinate at
every step, and has closing pair $\{u_{d-1}\},\varnothing$.  Equation
\eqref{eq:free-ledger} makes the dimension positive.  Decoding gives the
asserted refinement edges and the bijection gives distinctness and coverage.
\end{proof}

The particular Gray edges needed as ports are recorded separately, since
they are stronger than mere cyclicity.

\begin{lemma}[Low-coordinate Gray edges]\label{lem:gray-ports}
Assume $n$ is even, so $d\geq1$.  For the reflected Gray list on
$F(Q)=\{u_0<\cdots<u_{d-1}\}$:
\begin{enumerate}
\item if $d\geq1$ and $H\subseteq F(Q)\setminus\{u_0\}$, then $H$ and
      $H\cup\{u_0\}$ are consecutive;
\item if $d\geq2$ and $H\subseteq F(Q)\setminus\{u_0,u_1\}$, then
      $H\cup\{u_0\}$ and $H\cup\{u_0,u_1\}$ are consecutive; and
\item the closing edge is $\{u_{d-1}\},\varnothing$.
\end{enumerate}
\end{lemma}

\begin{proof}
For the first assertion, the case $d=1$ is the list
$\varnothing,\{u_0\}$.  For $d>1$, split $G_d$ into its two recursive halves.
If $u_{d-1}\notin H$, the required pair occurs in the first copy of
$G_{d-1}$ by induction.  If $u_{d-1}\in H$, remove $u_{d-1}$ from both masks;
the resulting pair is consecutive in $G_{d-1}$, and adding $u_{d-1}$ and
reversing the copy preserves adjacency.  The second assertion has the same
argument, with the explicit base list $G_2$ and the pair involving $u_1$.
Finally, the recursion starts at $\varnothing$, while its last mask is the
singleton containing the newest, hence greatest, coordinate $u_{d-1}$.
\end{proof}

If $Q\neq\varnothing$, let $p(Q)$ be obtained by deleting the
lexicographically first chord of $Q$.

\begin{lemma}[Canonical block tree]\label{lem:canonical-block-tree}
The map $p$ defines a rooted tree on canonical matchings, with root the empty
matching.  Each parent step decreases chord count by one and increases free
dimension by two.  Every canonical matching lies in the root subtree.
\end{lemma}

\begin{proof}
Deleting a chord preserves endpoint disjointness and noncrossingness.
It decreases $|Q|$ by one, so iteration terminates at the unique chordless
matching.  The parent is uniquely determined by the lexicographically first
chord.  Conversely, every nonroot matching has exactly that parent.  Thus the
directed parent relation is acyclic, connected to the empty matching, and has
unique parent at every nonroot vertex; it is a rooted tree.  Formula
\eqref{eq:free-ledger} gives the dimension change.
\end{proof}

Thus $C$ is a child of $Q$ precisely when
\[
 C=Q\cup\{e\},
 \qquad e\text{ is the lexicographically first chord of }C;
\]
equivalently, deleting the first chord of $C$ gives $Q$.  In particular, not
every admissible chord insertion into $Q$ is a child: the inserted chord must
precede every old chord of $Q$.

Let $C$ be a child of $Q$, with its first chord inserted between the $i$th and
$j$th free points of $Q$, $i<j$.  The reflected Gray cycles admit the
following explicit opposite edges.  Indices in a mask refer to the increasing
free-point order of the corresponding block.
\begin{equation}\label{eq:ports}
\begin{array}{c|c|c}
 (i,j)&\text{parent edge in }Q&\text{child edge in }C\\ \hline
 1\leq i<j&\{i,j\},\{0,i,j\}&\varnothing,\{0\}\\
 i=0,\ j\geq2&\{0,j\},\{0,1,j\}&\varnothing,\{0\}\\
 (i,j)=(0,1)&\varnothing,\{d(Q)-1\}&
      \varnothing,\{d(C)-1\}.
\end{array}
\end{equation}

\begin{lemma}[Explicit joining diamond]\label{lem:joining-diamond}
For every parent--child pair $Q,C$, the two displayed Gray edges in the
appropriate row of~\eqref{eq:ports}, together with two cross-block
refinement edges, form a four-cycle in $\NCR(n)$.
\end{lemma}

\begin{proof}
Write the free points of the parent as
$u_0<u_1<\cdots<u_{d(Q)-1}$, and let the deleted chord be
$(u_i,u_j)$.  The free points of the child are the remaining $u_r$'s, in the
same order.  If $B$ is a child mask, let $\widehat B$ be its order-preserving
lift to the parent coordinates.  Except in the special third row, the parent
mask paired with $B$ is
\begin{equation}\label{eq:parent-mask}
 M(B)=\widehat B\cup\{i,j\}.
\end{equation}
The decoder gives a direct comparison.  First, $u_i$ and $u_j$ have the same
anchor in the parent.  Indeed, an old chord contains $u_i$ in its open
interval if and only if it contains $u_j$: if it contained exactly one of
them, its endpoints would alternate with the new chord $(u_i,u_j)$, contrary
to the assumption that the child matching is noncrossing.  Their chains of
old containing chords are therefore identical.  Write their common
innermost anchor as
\[
 r=a_Q(u_i)=a_Q(u_j).
\]
In the child, $u_i$ and $u_j$ are endpoints of one matching chord and decode
to $u_i$; in the parent, selecting both in $M(B)$ sends both to $r$.

Now consider $x\notin\{u_i,u_j\}$.  Chords other than $(u_i,u_j)$ occur in
both matchings.  If $x$ lies outside the deleted chord, its chain of
containing chords is unchanged.  If $x$ lies inside it, then $(u_i,u_j)$ can
be its innermost chord only when its child key is $u_i$; otherwise a strictly
nested chord remains innermost in both matchings.  The order-preserving lift
keeps the selected/unselected status of every surviving free point.  Thus
deleting the chord performs exactly one relabelling of keys:
\[
 \kappa_{Q,M(B)}(x)=
 \begin{cases}
  r,&\kappa_{C,B}(x)=u_i,\\
  \kappa_{C,B}(x),&\kappa_{C,B}(x)\neq u_i.
 \end{cases}
\]
Here $r<u_i$, so the child fibers with keys $r$ and $u_i$ are distinct.  The
formula shows in both directions that every other key class is unchanged.
Thus $P_Q(M(B))$ is obtained from $P_C(B)$ by replacing precisely those two
fibers by their union, and the pair is joined by one cover edge.

We now verify that the four masks in each row of~\eqref{eq:ports} are the
ones just described and that the horizontal pairs are Gray edges.

If $i>0$, the least child free point lifts to $u_0$.  Hence
\[
 M(\varnothing)=\{i,j\},\qquad
 M(\{0\})=\{0,i,j\}.
\]
By part~1 of Lemma~\ref{lem:gray-ports}, both the parent pair and the child
pair $\varnothing,\{0\}$ are edges of their fixed Gray cycles.  The cross
covers are
\[
 P_C(\varnothing)\lessdot P_Q(\{i,j\}),\qquad
 P_C(\{0\})\lessdot P_Q(\{0,i,j\});
\]
these are the first row.

If $i=0$ and $j\geq2$, the least child free point lifts to $u_1$.  Thus
\[
 M(\varnothing)=\{0,j\},\qquad
 M(\{0\})=\{0,1,j\}.
\]
By part~2 of Lemma~\ref{lem:gray-ports}, the parent pair is a Gray-cycle edge;
by part~1, so is the child pair.  The cross covers are
\[
 P_C(\varnothing)\lessdot P_Q(\{0,j\}),\qquad
 P_C(\{0\})\lessdot P_Q(\{0,1,j\});
\]
these are the second row.

Finally suppose $(i,j)=(0,1)$.  Here the two selected horizontal edges are
the closing Gray edges.  The rightmost child free point lifts to the
rightmost parent free point.  At the empty mask the endpoints $u_0,u_1$ are
singletons in the parent and one chord fiber in the child, so inserting the
chord performs exactly their merge.  At the rightmost singleton mask, the
selected rightmost point has the same outer anchor before and after the
insertion because it lies to the right of $u_1$; the only changed keys are
again those of $u_0,u_1$.  Hence the cross covers are
\[
 P_Q(\varnothing)\lessdot P_C(\varnothing),\qquad
 P_Q(\{d(Q)-1\})\lessdot P_C(\{d(C)-1\}).
\]
By part~3 of Lemma~\ref{lem:gray-ports}, the two horizontal pairs are the
closing edges in their respective block cycles.  In all three cases the four
vertices are distinct across the two Boolean blocks, both horizontal pairs
are fixed block-cycle edges, and both vertical pairs are single merges.  They
therefore form the required four-cycle.
\end{proof}

Represent a Gray edge by a \emph{port key}, either
$(q,H)$, meaning ``flip coordinate $q$ above ambient lower mask $H$'', or
\emph{closing}.  Here $H$ is the endpoint not containing $q$, so the other
endpoint is $H\cup\{q\}$.  The three rows of~\eqref{eq:ports} give one
incoming key in the parent and one outgoing key in the child.

\begin{lemma}[Port injection]\label{lem:port-injection}
For a fixed block $Q$, distinct children are assigned to distinct directed
edge occurrences of the Gray cycle of $Q$.  If $Q$ is nonroot, its outgoing
edge occurrence toward $p(Q)$ is not assigned to any child.
\end{lemma}

\begin{proof}
Record a nonclosing directed Gray occurrence by a pair $(q,H)$, where $q$ is
the flipped ambient point and $H$ is the lower ambient mask.  Equality of two
such keys forces equality of both $q$ and $H$; the closing occurrence is a
separate key.

Let $C$ be a child of $Q$, with first chord $e=(u_i,u_j)$.  The parent-side
key supplied by~\eqref{eq:ports} is
\[
 \begin{cases}
  (u_0,\{u_i,u_j\}),&i>0,\\
  (u_1,\{u_0,u_j\}),&i=0,\ j\geq2,\\
  \text{closing},&(i,j)=(0,1).
 \end{cases}
\]
In each of the first two cases the lower mask is exactly the two-element
endpoint set of $e$.  Because the endpoints are ordered increasingly, that
set uniquely recovers the ordered chord.  The child itself is then recovered
as $Q\cup\{e\}$.  Thus two equal coordinate keys come from the same child.
If both keys are closing, both first chords join the first two free points of
$Q$, so again the chords and children coincide.  A coordinate key cannot
equal the closing key.  This proves injectivity.

It remains to check that the edge retained for the parent of $Q$ is not used
by a child of $Q$.  If the outgoing key is a coordinate key, its lower mask
is $\varnothing$, whereas every incoming coordinate key has lower mask
$\{u_i,u_j\}\neq\varnothing$; an incoming closing key has a different key
type.  If the outgoing key is closing, the first chord of $Q$ joins its first
two points that are free in $p(Q)$.  Write this chord as $e=(a,b)$.  If a
child of $Q$ existed, let $f=(c,d)$ be the chord whose deletion returns that
child to $Q$.
Both endpoints of $f$ are free in $Q$, and all such points lie to the right of
$e$, because $e$ used the first two free points of $p(Q)$.  Hence
$b<c$, so $e<_{\mathrm{lex}}f$.  On the other hand, $e$ remains a chord
of the child while $f$ must be its lexicographically first chord, forcing
$f\leq_{\mathrm{lex}}e$, a contradiction.  Thus in this special case $Q$ has
no children at all.  The outgoing occurrence is therefore unassigned in
every case.
\end{proof}

\paragraph{The case $n=4$.}
All three rows of \eqref{eq:ports} occur when $n=4$.  For the root matching
$Q=\varnothing$, the free points are
$u_0=1,u_1=2,u_2=3$, so its Boolean block is the $3$-cube with eight
vertices.  The only other canonical matchings are the three one-chord
matchings; each has one free point and hence a two-vertex Boolean block.  The
four blocks have sizes
\[
  2^3+2^1+2^1+2^1=14=\catalan_4.
\]
Their ports in the root Gray cycle are listed below.  Chord labels are
ambient points of $X_4$, whereas the sets in the last two columns are
mask-coordinate indices.
\[
\begin{array}{c|c|c|c}
\text{child chord}&(i,j)&\text{root edge}&\text{child edge}\\ \hline
(2,3)&(1,2)&\{1,2\},\{0,1,2\}&\varnothing,\{0\}\\
(1,3)&(0,2)&\{0,2\},\{0,1,2\}&\varnothing,\{0\}\\
(1,2)&(0,1)&\varnothing,\{2\}&\varnothing,\{0\}.
\end{array}
\]
The first two root edges share the vertex $\{0,1,2\}$, but they are distinct
directed edge occurrences.  The recursive construction below is indexed by
directed occurrences, not by pairwise vertex-disjoint edges: each
occurrence is replaced by one path piece, and adjacent pieces hand off at the
shared Gray vertex.  Splicing the three two-vertex child codes into these
three root occurrences produces one cyclic code on all fourteen vertices.
This is the $n=4$ instance of the general recursion.

\section{Recursive square switching}\label{sec:assembly}

We use the elementary square switch: if two disjoint cycles contain opposite
edges $ab$ and $cd$ of a four-cycle $a b d c a$, deleting $ab,cd$ and adding
$ac,bd$ produces one cycle on the union of their vertices.  Repeated switches
modify the parent cycle, so the local square picture alone does not specify
which parent edges remain.  The recursive operation is given by the following
list-concatenation principle.

At the recursive level, a \emph{cyclic code} is a cyclic list of distinct
vertices with adjacent consecutive entries.  A two-vertex code records the
two directed occurrences of its single undirected edge; codes on at least
three vertices are ordinary simple cycles.  The two-vertex convention is used
only inside the recursion.  The assembled root code below has at least three
vertices and therefore gives a simple Hamilton cycle.

\begin{lemma}[Cyclic piece concatenation]\label{lem:piece-concatenation}
Let $v_0,v_1,\ldots,v_{r-1}$ be cyclically indexed, and for each directed
occurrence $e_k=(v_k,v_{k+1})$ let $K_k$ be a finite vertex list.  Suppose
that
\begin{enumerate}
\item each $K_k$ is nonempty, has no repeated vertex, and starts at $v_k$;
\item consecutive vertices inside $K_k$ are adjacent;
\item the last vertex of $K_k$ is adjacent to $v_{k+1}$; and
\item the lists $K_k$ are pairwise vertex-disjoint.
\end{enumerate}
Then
\[
 K_0\mathbin{\Vert}K_1\mathbin{\Vert}\cdots
 \mathbin{\Vert}K_{r-1}
\]
is a cyclic code with support the disjoint union of the supports of the
$K_k$.
\end{lemma}

\begin{proof}
Within each piece, adjacency is assumption~2.  Between successive pieces,
the last vertex of $K_k$ is adjacent by assumption~3 to $v_{k+1}$, which is
the first vertex of $K_{k+1}$ by assumption~1; the same argument joins the
last piece back to $K_0$.  Each piece is internally repetition-free, and
assumption~4 prevents repetitions between pieces.  Flattening the pieces
therefore gives the asserted cyclic code and support.
\end{proof}

Directed occurrences are essential here.  If a child code
has only two vertices $[a,b]$, its cyclic occurrences are $(a,b)$ and
$(b,a)$.  Cutting at $(a,b)$ leaves the path list $[b,a]$ (and conversely for
the other orientation), so both vertices still occur exactly once.  We are
cutting a cyclic list at a specified occurrence; the underlying simple graph
still has only one edge.

\begin{lemma}[Block-tree assembly]\label{lem:tree-assembly}
For every canonical matching $Q$ and even $n$, there is a cyclic code whose
support is exactly the union of all Boolean blocks in the subtree rooted at
$Q$.
If $Q$ is nonroot, the cyclic code retains the outgoing port prescribed by the
diamond joining $Q$ to $p(Q)$.
\end{lemma}

\begin{proof}
We prove simultaneously the following three assertions for the subtree at
$Q$:
\begin{enumerate}
\item there is one cyclic code $Z_Q$, which is a simple cycle when its support
has at least three vertices;
\item its support is exactly the union of the Boolean blocks in the subtree;
\item if $Q$ is nonroot, the prescribed outgoing directed edge of the block
cycle of $Q$ is still an edge of $Z_Q$.
\end{enumerate}
Induct on $d(Q)=|F(Q)|$.  A child is obtained by adding one chord and hence
has dimension $d(Q)-2$, so this is a well-founded induction.  Assume the
three assertions for every child.

The minimum possible dimension is $d(Q)=1$.  Then no child exists, since a
new chord would require two free endpoints.  The two directed occurrences of
the one-dimensional block code are both unassigned, their pieces are the two
singleton source lists, and their concatenation is the required two-vertex
cyclic code.  This establishes the base case.

Write the cyclic Gray list of the block of $Q$ as
$v_0,v_1,\ldots,v_{r-1}$, with indices modulo $r$, and regard
$e_k=(v_k,v_{k+1})$ as a directed edge occurrence.  By
Lemma~\ref{lem:port-injection}, each child is assigned to a distinct $e_k$.
The assigned occurrences need not be vertex-disjoint, as the $n=4$ example
shows; distinctness is exactly what is required because the construction
builds one replacement piece for each source occurrence.
For an unassigned occurrence $e_k$, define the corresponding piece to be the
singleton list $[v_k]$.  For an occurrence assigned to a child $C$, the
inductive cyclic code $Z_C$ still contains its prescribed outgoing edge $ab$.
The joining diamond has opposite directed occurrences $e_k$ and $ab$.
Cutting $Z_C$ at $ab$ turns it into a path list through every vertex of the
child subtree, in one of its two orientations.  Orient it as
$W_C=[w_0,\ldots,w_s]$ so that the diamond supplies the cross edges
$v_kw_0$ and $w_sv_{k+1}$.  The assigned piece is the list
\[
 K_k=[v_k]\mathbin{\Vert}W_C.
\]
Its first vertex is $v_k$, its internal adjacencies are the first cross edge
followed by the cut child path, and its last vertex is adjacent through the
second cross edge to $v_{k+1}$.  Thus every piece associated with $e_k$
starts at the source of $e_k$ and hands off to the source of $e_{k+1}$; the
preceding paragraph covers the two-vertex child case as well.

We verify that concatenation introduces neither omissions nor repetitions.
The vertices $v_k$ are distinct because the block Gray list is a cyclic
listing.  A child-subtree piece contains no vertex in the block of $Q$:
the canonical matching of every one of its vertices lies at or below that
child, hence strictly below $Q$ in the matching tree, whereas the vertices
$v_k$ have canonical
matching exactly $Q$.  Two different child-subtree pieces are disjoint
because a rooted tree has disjoint subtrees below distinct children.
Finally, port injectivity ensures that no child subtree is inserted twice.
Therefore all pieces are nonempty, internally simple, and pairwise disjoint.

Their internal edges are either edges of a cut child code or the first cross
edge in its joining diamond.  The exit of the $k$th piece is adjacent through
the second cross edge to $v_{k+1}$, the head of the next piece; the final
piece hands off to $v_0$ in exactly the same way.  The four hypotheses of
Lemma~\ref{lem:piece-concatenation} are now verified.  Its concatenation is a
single cyclic code, and a simple cycle whenever it has at least three
vertices.

Every vertex in the block of $Q$ occurs as exactly one $v_k$.  Every proper
descendant belongs to the subtree of a unique child, and that entire child
subtree occurs in the unique piece assigned to that child.  The support is
therefore exactly the subtree of $Q$.  If $Q$ is nonroot, its outgoing
occurrence is unassigned by Lemma~\ref{lem:port-injection}; its piece is the
singleton $[v_k]$, so the handoff edge $v_kv_{k+1}$ remains present in the
flattened cycle.  This proves the third assertion and closes the induction.
\end{proof}

\begin{theorem}[Uniform even-order construction]\label{thm:uniform-even}
For every even $n\geq4$, $\NCR(n)$ has a Hamilton cycle.
\end{theorem}

\begin{proof}
Apply Lemma~\ref{lem:tree-assembly} to the empty matching.  By
Lemma~\ref{lem:canonical-block-tree}, every canonical matching lies in its
subtree.  By Theorem~\ref{thm:boolean-decomposition}, the corresponding
Boolean blocks partition all of $\NC(n)$.  The assembled root cycle is
therefore spanning.  Since $n\geq4$, $|\NC(n)|=\catalan_n\geq3$~\cite{S}, so the
closed spanning walk is a Hamilton cycle in the simple graph.
\end{proof}

\section{Proof of the classification}

\begin{proof}[Proof of Theorem~\ref{thm:classification}]
For even $n\geq4$, apply Theorem~\ref{thm:uniform-even}.  For odd
$n\geq3$, apply Theorem~\ref{thm:odd-obstruction}.

It remains to inspect $n<4$.  There is one noncrossing partition for $n=0$
and one for $n=1$, so the convention in the theorem regards both singleton
graphs as Hamiltonian.  For $n=2$, the refinement graph consists of one edge
on two vertices and has no simple Hamilton cycle.  For $n=3$, the odd
obstruction applies (equivalently, the five-vertex bipartition is unbalanced).
These cases give exactly the stated disjunction.
\end{proof}

\section{Concluding remarks}

For odd orders, the two block-parity classes have different sizes.  For even
orders, every Boolean block has odd dimension, and the lexicographic chord
tree provides injectively assigned joining ports.  The even-order construction
does not require a Hamilton path with prescribed endpoints inside each
Boolean block.  It reserves specified Gray edges and joins parent and child
cycles through four-cycles in the refinement graph.

\section{Machine-checked correspondence}\label{sec:formal}

The main Lean~4 import is \texttt{Hamilton.HamiltonianCycles}.  The final
cycle and path classifications are
\leanname{NCRefinementGraph_fin_isHamiltonian_iff} and
\leanname{NCRefinementGraph_fin_isHamiltonianPath_iff}.  To trace an
intermediate result, locate its paper label or proof component in
Table~\ref{tab:lean-map} and then search for the declarations in the second
column.  The table is a navigation index; the mathematical identification of
the paper's definitions with their Lean counterparts is explained in the
surrounding prose.

The first identification to check is graph adjacency.  The declaration
\begin{center}
\leanname{NCRefinementGraph_adj_iff_single_merge}
\end{center}
expands adjacency to a two-block merge in either direction.
The declarations \leanname{mergesTo_numBlocks} and \leanname{mergesTo_le}
state that such a merge removes one block and moves upward in the refinement
order.  Conversely, the graded-lattice argument in the introduction shows
that a cover cannot skip a rank and therefore must merge exactly two blocks.
Thus the Lean adjacency is the cover-graph adjacency used in the paper.

\begingroup\footnotesize
\begin{longtable}{p{0.34\textwidth}p{0.55\textwidth}}
\caption{Paper-to-Lean correspondence.}\label{tab:lean-map}\\
\toprule
Paper statement or proof step & Lean declarations\\
\midrule
\endfirsthead
\toprule
Paper statement or proof step & Lean declarations\\
\midrule
\endhead
\leanname{def:refinement-graph} &
\leanname{NCRefinementGraph_adj_iff_single_merge},
\leanname{mergesTo_numBlocks}, \leanname{mergesTo_le}\\
\leanname{lem:bipartite} &
\leanname{NCRefinementGraph_isBipartite}\\
\leanname{thm:signed-count} &
\leanname{signedDyckSumTree_closed_form},
\leanname{signedNCCount_univ_eq_signedDyckSumTree}\\
\leanname{thm:odd-obstruction} &
\leanname{NCRefinementGraph_fin_odd_geq3_not_isHamiltonian}\\
\leanname{lem:key-fibers} &
\leanname{commonKey_root_or_chord}, \leanname{point_in_chord_hull},
\leanname{point_avoids_nested_chord}, \leanname{root_point_avoids_chord},
\leanname{decodedFinpartition_isNoncrossing}\\
\leanname{thm:boolean-decomposition} &
\leanname{decodeState_bijective}, \leanname{stateEquivNC},
\leanname{decodedNC_mergesTo_insert}\\
\leanname{lem:block-gray-cycle} &
\leanname{rankGrayCycle}, \leanname{rankGrayVertices_nodup},
\leanname{rankGrayVertices_chain}\\
\leanname{lem:gray-ports} &
\leanname{coordinateZero_rankGray_consecutive},
\leanname{coordinateOne_rankGray_consecutive},
\leanname{rankGray_closing_flip}\\
\leanname{lem:canonical-block-tree} &
\leanname{parent_card}, \leanname{parentOption_eq_none_iff},
\leanname{inSubtree_root}\\
\leanname{lem:joining-diamond} &
\leanname{parent_decodeKey_eq_relabel_child},
\leanname{decodedNC_child_adj_parentMask},
\leanname{joiningDiamond_nonempty}, \leanname{joiningDiamond}\\
\leanname{lem:port-injection} &
\leanname{portAssignedToPair_exists}, \leanname{assignedPair_injective},
\leanname{outgoingPair_not_assigned}\\
\leanname{lem:piece-concatenation} &
\leanname{EdgePieceSystem.toCycleCode}\\
\leanname{lem:tree-assembly} &
\leanname{edgePieces_pairwise_disjoint},
\leanname{subtreeEdgePieceSystem}, \leanname{subtreeCycle},
\leanname{assembledSubtreeCycleCode_support},
\leanname{assembledSubtreeCycleCode_outgoing},
\leanname{rootCycleCode_support}\\
\leanname{thm:uniform-even} &
\leanname{NCRefinementGraph_fin_even_geq4_isHamiltonian_petersen}\\
\leanname{thm:classification} &
\leanname{NCRefinementGraph_fin_isHamiltonian_iff}\\
\leanname{cor:hamilton-paths} ($n\leq3$) &
\leanname{NCRefinementGraph_fin_zero_isHamiltonianPath},
\leanname{NCRefinementGraph_fin_one_isHamiltonianPath},
\leanname{NCRefinementGraph_fin_two_isHamiltonianPath},
\leanname{NCRefinementGraph_fin_three_isHamiltonianPath},
\leanname{cardThree_hamiltonianPath_between_twoBlock}\\
\leanname{cor:hamilton-paths}: even case &
\leanname{IsHamiltonian.isHamiltonianPath}\\
\leanname{cor:hamilton-paths}: odd case &
\leanname{NCRefinementGraph_evenBlocks_oddBlocks_diff_at_most_one_of_isHamiltonianPath},
\leanname{NCRefinementGraph_fin_odd_geq5_not_isHamiltonianPath}\\
\leanname{cor:hamilton-paths}: final &
\leanname{NCRefinementGraph_fin_isHamiltonianPath_iff}\\
\bottomrule
\end{longtable}
\endgroup

From the \texttt{lean4} directory, run
\begin{quote}\scriptsize\ttfamily
lake build Hamilton.HamiltonianCycles\\
lake env lean --trust=0 Hamilton/HamiltonianCyclesTheoremMap.lean\\
lake env lean --trust=0 Hamilton/HamiltonianCyclesAxiomAudit.lean
\end{quote}
The first command builds the formalization, the second checks the declarations
in Table~\ref{tab:lean-map}, and the third reports for the audited construction,
obstruction, cycle-classification, and path-classification endpoints exactly
\[
 \texttt{[propext, Classical.choice, Quot.sound]}.
\]
These are standard Lean and Mathlib logical dependencies.  No additional
axiom, \texttt{sorry}, \texttt{admit}, or \texttt{native\_decide} occurs in
the dependency closure of these theorems.

The audit checks the Lean theorems under their formal definitions and the
foundational dependencies just listed.  The correspondence between those
definitions and the mathematical objects in the paper is part of the
mathematical exposition.

The comparison with the refinement Gray-code problem of Merino, Namrata, and
Williams is a literature-level identification, not an additional formal
theorem.  Lean formalizes the merge/split adjacency and the Hamilton-path
classification; the identification with their externally stated list problem
uses the definition in the cited work.

\ifnum\blindmode=0
Archive concept DOI:
\href{https://doi.org/10.5281/zenodo.21316750}{\nolinkurl{10.5281/zenodo.21316750}}.
It resolves to the latest version.
\else
The kernel-checked formalization is publicly archived.  The archive identifier
is omitted from this anonymous copy.
\fi

\section*{Declaration of Generative AI and AI-Assisted Technologies in the Manuscript Preparation Process}

During the preparation of this work, the author used OpenAI Codex for
exploratory proof development, drafting and copy-editing, code development,
and the Lean~4 formalization.  Claude (Anthropic) was used during the initial
exploration.  The author checked the mathematical claims and citations,
compared the manuscript statements with the Lean declarations, and reran the
formal build and axiom audit.  AI output is not treated as mathematical
evidence.  The author is responsible for the definitions, proofs, citations,
code, and manuscript.

\ifnum\blindmode=0
\subsection*{Acknowledgements}
The author thanks the developers of Lean and Mathlib for the formalization
environment.
\fi

\begin{thebibliography}{20}

\bibitem{HHMW}
E.~Hartung, H.~P.~Hoang, T.~M\"utze, and A.~Williams,
Combinatorial generation via permutation languages. I. Fundamentals,
\emph{Trans. Amer. Math. Soc.} \textbf{375} (2022), 2255--2291;
\doi{10.1090/tran/8199}; \arxiv{1906.06069}.

\bibitem{HHNP}
C.~Huemer, F.~Hurtado, M.~Noy, and E.~Oma\~na-Pulido,
Gray codes for non-crossing partitions and dissections of a convex polygon,
\emph{Discrete Appl. Math.} \textbf{157} (2009), 1509--1520;
\doi{10.1016/j.dam.2008.06.018}.

\bibitem{HM}
H.~P.~Hoang and T.~M\"utze,
Combinatorial generation via permutation languages. II. Lattice congruences,
\emph{Israel J. Math.} \textbf{244} (2021), 359--417;
\doi{10.1007/s11856-021-2186-1}; \arxiv{1911.12078}.

\bibitem{K}
G.~Kreweras,
Sur les partitions non crois\'ees d'un cycle,
\emph{Discrete Math.} \textbf{1} (1972), 333--350;
\doi{10.1016/0012-365X(72)90041-6}.

\bibitem{MNW}
A.~Merino, Namrata, and A.~Williams,
On the hardness of Gray code problems for combinatorial objects,
preprint, 2024; \arxiv{2401.14963}.

\bibitem{M}
T.~M\"utze,
Combinatorial Gray codes---an updated survey,
\emph{Electron. J. Combin.} \textbf{30}(3) (2023), DS26;
\doi{10.37236/11023}; \arxiv{2202.01280}.

\bibitem{Pet}
T.~K.~Petersen,
On the shard intersection order of a Coxeter group,
\emph{SIAM J. Discrete Math.} \textbf{27} (2013), 1880--1912;
\doi{10.1137/110847202}; \arxiv{1108.5761}.

\bibitem{S}
R.~Simion,
Noncrossing partitions,
\emph{Discrete Math.} \textbf{217} (2000), 367--409;
\doi{10.1016/S0012-365X(99)00273-3}.

\bibitem{SU}
R.~Simion and D.~Ullman,
On the structure of the lattice of noncrossing partitions,
\emph{Discrete Math.} \textbf{98} (1991), 193--206;
\doi{10.1016/0012-365X(91)90376-D}.

\end{thebibliography}

\end{document}
